\documentclass[11pt,a4paper]{article}

\usepackage[a4paper,margin=1in]{geometry}
\usepackage{iftex}
\ifPDFTeX
  \usepackage[T5]{fontenc}
  \usepackage[utf8]{inputenc}
  \usepackage[vietnamese]{babel}
  \usepackage{lmodern}
\else
  \usepackage{fontspec}
  \usepackage{polyglossia}
  \setmainlanguage{vietnamese}
  \setotherlanguage{english}
  \setmainfont{Times New Roman}
  \setsansfont{Arial}
  \setmonofont{Consolas}
\fi

\usepackage{xcolor}
\usepackage{booktabs}
\usepackage{array}
\usepackage{colortbl}
\usepackage{longtable}
\usepackage{enumitem}
\usepackage{hyperref}
\usepackage{graphicx}
\usepackage{tikz}
\usetikzlibrary{arrows.meta,positioning,fit,calc,shapes.geometric,shapes.misc,matrix}
\usepackage{caption}
\usepackage{float}
\usepackage{amsmath}

\hypersetup{
  colorlinks=true,
  linkcolor=blue!55!black,
  urlcolor=blue!55!black,
  citecolor=blue!55!black
}
\setlength{\emergencystretch}{3em}

\definecolor{CoreBlue}{HTML}{1565C0}
\definecolor{CoreGreen}{HTML}{2E7D32}
\definecolor{CoreOrange}{HTML}{E65100}
\definecolor{CoreRed}{HTML}{C62828}
\definecolor{SoftGray}{HTML}{F5F7FA}
\definecolor{LineGray}{HTML}{6B7280}
\definecolor{SoftYellow}{HTML}{FFF7D6}

\tikzset{
  box/.style={
    rectangle,
    rounded corners=2pt,
    draw=LineGray,
    fill=SoftGray,
    align=center,
    minimum height=9mm,
    text width=33mm,
    font=\small
  },
  widebox/.style={
    rectangle,
    rounded corners=2pt,
    draw=LineGray,
    fill=SoftGray,
    align=center,
    minimum height=10mm,
    text width=52mm,
    font=\small
  },
  smallbox/.style={
    rectangle,
    rounded corners=2pt,
    draw=LineGray,
    fill=SoftGray,
    align=center,
    minimum height=7mm,
    text width=25mm,
    font=\footnotesize
  },
  good/.style={box, draw=CoreGreen, fill=CoreGreen!8},
  warn/.style={box, draw=CoreOrange, fill=CoreOrange!9},
  bad/.style={box, draw=CoreRed, fill=CoreRed!8},
  crypto/.style={box, draw=CoreBlue, fill=CoreBlue!8},
  highlight/.style={box, draw=CoreOrange, fill=SoftYellow},
  arrow/.style={-{Latex[length=2.4mm]}, thick, draw=LineGray},
  dashedarrow/.style={-{Latex[length=2.4mm]}, thick, dashed, draw=LineGray},
  smallnote/.style={font=\footnotesize, align=center, text=LineGray},
  layer/.style={rectangle, rounded corners=3pt, draw=LineGray, dashed, inner sep=5pt},
  coeff/.style={rectangle, draw=LineGray, minimum width=8mm, minimum height=7mm, align=center, font=\small},
  coeffgood/.style={coeff, fill=CoreGreen!12, draw=CoreGreen},
  coeffwarn/.style={coeff, fill=SoftYellow, draw=CoreOrange}
}

\title{\textbf{Giải Thích Chi Tiết Cách Hệ Thống Xử Lý Video Và Nhúng Tin}\\
\large H.264/CAVLC Patchability-Aware Video Steganography}
\author{Tài liệu kỹ thuật cho dự án}
\date{\today}

\begin{document}
\maketitle

\section{Câu Trả Lời Ngắn Gọn Trước}

Nếu nói thật ngắn, hệ thống làm việc như sau:

\begin{center}
\textit{Nó không sửa pixel của video. Nó sửa một vài dấu của các hệ số nén
trong H.264, cụ thể là dấu của các trailing-one residual coefficients trong
các block CAVLC của IDR/I-frame.}
\end{center}

Một file H.264 được hệ thống nhìn như một cấu trúc nhiều tầng:

\begin{center}
\texttt{file .h264}
\(\rightarrow\)
\texttt{NAL units}
\(\rightarrow\)
\texttt{IDR slices}
\(\rightarrow\)
\texttt{macroblocks}
\(\rightarrow\)
\texttt{4x4 residual blocks}
\(\rightarrow\)
\texttt{coefficients}
\end{center}

Thông tin được nhúng vào tầng cuối cùng: \textbf{coefficient}. Nhưng không phải
coefficient nào cũng được sửa. Hệ thống chỉ sửa các vị trí đã qua nhiều lớp
lọc: an toàn về CAVLC, có thể vá lại bitstream, không làm chất lượng giảm quá
mức, và đọc lại được sau khi tái tạo video.

\section{Codec, H.264 Và Vì Sao Hệ Thống Không Sửa Pixel}

\subsection{Codec là gì?}

Codec là bộ mã hóa/giải mã video. Tên này đến từ:

\begin{center}
\texttt{coder + decoder = codec}
\end{center}

Video thô rất lớn. Codec như H.264 nén video thành bitstream nhỏ hơn để lưu,
truyền, hoặc phát lại. Khi phát video, decoder đọc bitstream đó và dựng lại
các frame.

\subsection{Sửa pixel và sửa bitstream khác nhau thế nào?}

Một hệ đơn giản có thể làm như sau:

\begin{center}
\texttt{decode video -> sửa pixel -> encode lại video}
\end{center}

Đó là cách làm ở \textbf{pixel domain}. Nó dễ hiểu nhưng có một vấn đề lớn:
khi encode lại, codec có thể thay đổi rất nhiều thứ. Bit nhúng có thể mất,
artifact có thể xuất hiện, và khó kiểm soát chính xác bitstream cuối cùng.

Hệ thống này đi theo hướng khác:

\begin{center}
\texttt{parse H.264 bitstream -> sửa coefficient đã nén -> patch lại H.264}
\end{center}

Đây là cách làm ở \textbf{compressed domain}. Nó gần codec hơn, vì hệ thống
hiểu một phần cấu trúc H.264 và sửa đúng những chỗ mà codec đang dùng để biểu
diễn residual.

\begin{figure}[H]
\centering
\resizebox{\textwidth}{!}{%
\begin{tikzpicture}[node distance=10mm and 12mm]
  \node[widebox] (pixel0) {Cách pixel-domain\\Decode thành ảnh};
  \node[widebox, right=of pixel0] (pixel1) {Sửa pixel\\LSB hoặc màu};
  \node[widebox, right=of pixel1] (pixel2) {Encode lại\\codec có thể đổi nhiều thứ};

  \node[widebox, below=18mm of pixel0] (comp0) {Cách của hệ thống\\Parse H.264};
  \node[widebox, right=of comp0] (comp1) {Sửa residual coefficient\\trong CAVLC block};
  \node[widebox, right=of comp1] (comp2) {Patch/reconstruct\\giữ bitstream hợp lệ};

  \draw[arrow] (pixel0) -- (pixel1);
  \draw[arrow] (pixel1) -- (pixel2);
  \draw[arrow] (comp0) -- (comp1);
  \draw[arrow] (comp1) -- (comp2);
\end{tikzpicture}
}
\caption{Khác biệt giữa sửa pixel và sửa dữ liệu nén trong H.264.}
\label{fig:pixel-vs-compressed}
\end{figure}

\section{File H.264 Được Hệ Thống Bóc Ra Như Thế Nào?}

Phần này là lõi để nhìn rõ hệ thống. Khi nhận một file \texttt{video.h264}, hệ
thống không nhìn nó như một chuỗi byte vô nghĩa. Nó bóc file theo các tầng sau.

\begin{figure}[H]
\centering
\begin{tikzpicture}[node distance=8mm]
  \node[widebox] (bitstream) {File H.264 bitstream\\\texttt{video.h264}};
  \node[widebox, below=of bitstream] (nal) {NAL units\\đơn vị dữ liệu của H.264};
  \node[widebox, below=of nal] (slice) {IDR slice / I-frame\\frame ít phụ thuộc frame khác};
  \node[widebox, below=of slice] (mb) {Macroblock\\vùng 16x16 pixel};
  \node[widebox, below=of mb] (block) {Luma 4x4 residual block\\block nhỏ bên trong macroblock};
  \node[widebox, below=of block] (coeff) {Residual coefficients\\các số đã biến đổi/lượng tử hóa};
  \node[good, below=of coeff] (target) {T1 sign positions\\vị trí nhúng bit};

  \draw[arrow] (bitstream) -- (nal);
  \draw[arrow] (nal) -- (slice);
  \draw[arrow] (slice) -- (mb);
  \draw[arrow] (mb) -- (block);
  \draw[arrow] (block) -- (coeff);
  \draw[arrow] (coeff) -- (target);
\end{tikzpicture}
\caption{Tầng xử lý từ file H.264 đến vị trí coefficient có thể nhúng.}
\label{fig:h264-layers}
\end{figure}

Một vị trí nhúng trong hệ thống có thể ghi gọn là:

\begin{center}
\texttt{(macroblock\_id, block\_id, coefficient\_id)}
\end{center}

Ví dụ:

\begin{center}
\texttt{(496, 7, 1)}
\end{center}

Nghĩa là:

\begin{itemize}[leftmargin=*]
  \item macroblock số 496,
  \item block 4x4 số 7 trong macroblock đó,
  \item coefficient số 1 trong block đó.
\end{itemize}

\section{Ví Dụ Đầy Đủ: Một Video H.264 Được Parse Ra Gì?}

Phần này dùng một ví dụ minh họa gần với cách hệ thống làm việc. Đây không phải
là dump đầy đủ của một file benchmark cụ thể, mà là một mẫu dữ liệu rút gọn để
nhìn rõ pipeline. Trong code thật, các bước tương ứng nằm ở parser H.264,
CAVLC parser, safety filter và bitstream reconstructor.

\subsection{File H.264 ban đầu trông như thế nào?}

Một file \texttt{.h264} thô là một chuỗi byte. Trong H.264 Annex B, các NAL unit
thường bắt đầu bằng start code:

\begin{center}
\texttt{00 00 00 01}
\end{center}

Ví dụ một đoạn bitstream rút gọn có thể nhìn như sau:

\begin{center}
\begin{tabular}{p{0.22\linewidth}p{0.68\linewidth}}
\toprule
Offset byte & Dữ liệu minh họa \\
\midrule
0x00000000 &
\texttt{00 00 00 01 67 ...} \\
0x0000001A &
\texttt{00 00 00 01 68 ...} \\
0x00000024 &
\texttt{00 00 00 01 65 B8 21 ...} \\
0x00004F10 &
\texttt{00 00 00 01 65 88 40 ...} \\
\bottomrule
\end{tabular}
\end{center}

Hệ thống đọc các start code này để tách file thành các NAL unit. Sau đó nó đọc
header của mỗi NAL để biết NAL đó là SPS, PPS, IDR slice hay loại khác.

\subsection{Kết quả parse NAL mẫu}

Sau bước parse NAL, hệ thống có thể có một bảng nội bộ giống như sau:

\begin{center}
\begin{tabular}{rcllp{45mm}}
\toprule
NAL & Offset & Type & Vai trò & Hệ thống làm gì? \\
\midrule
0 & 0x00000000 & 7 & SPS & Đọc kích thước, profile, thông tin frame. \\
1 & 0x0000001A & 8 & PPS & Đọc entropy mode, CAVLC/CABAC, tham số slice. \\
2 & 0x00000024 & 5 & IDR slice & Có thể đi vào macroblock và tìm vị trí nhúng. \\
3 & 0x00004F10 & 5 & IDR slice & Tiếp tục tìm block/coefficient dùng được. \\
\bottomrule
\end{tabular}
\end{center}

Với baseline hiện tại, hệ thống quan tâm nhất tới NAL type 5, tức IDR slice.
Các NAL SPS/PPS cần thiết để hiểu stream, nhưng không phải nơi nhúng payload.

\subsection{Từ IDR slice đến macroblock}

Một IDR slice chứa nhiều macroblock. Với video CIF 352x288, mỗi frame có:

\begin{center}
\texttt{22 macroblock theo chiều ngang} \(\times\)
\texttt{18 macroblock theo chiều dọc}
\(= 396\) macroblock/frame.
\end{center}

Một bảng macroblock rút gọn có thể như sau:

\begin{center}
\begin{tabular}{rrrrp{46mm}}
\toprule
Frame & Macroblock & Hàng/cột MB & Slice-local MB & Ghi chú \\
\midrule
12 & 4752 & row 0, col 0 & 0 & IDR frame bắt đầu. \\
12 & 4753 & row 0, col 1 & 1 & Có luma residual blocks. \\
12 & 4960 & row 9, col 10 & 208 & Có nhiều T1 candidates. \\
12 & 5147 & row 17, col 21 & 395 & Macroblock cuối frame. \\
\bottomrule
\end{tabular}
\end{center}

Trong code, chỉ số macroblock thường được dùng theo dạng global index để các
frame khác nhau không bị trùng vị trí.

\begin{figure}[H]
\centering
\resizebox{\textwidth}{!}{%
\begin{tikzpicture}
  \draw[step=0.45cm, LineGray, thin] (0,0) grid (9.9,8.1);
  \draw[CoreBlue, very thick] (4.5,3.6) rectangle (4.95,4.05);
  \node[smallnote] at (4.72,4.35) {MB row 9, col 10};
  \node[smallnote] at (4.95,-0.4) {22 macroblock mỗi hàng};
  \node[smallnote, rotate=90] at (-0.45,4.05) {18 hàng macroblock};
\end{tikzpicture}
}
\caption{Một frame CIF 352x288 được chia thành lưới 22x18 macroblock. Ô được tô là một macroblock minh họa.}
\label{fig:macroblock-grid}
\end{figure}

\subsection{Từ macroblock đến các block 4x4}

Mỗi macroblock luma có thể được nhìn như 16 block 4x4. Hệ thống đánh số các
block con này và đọc residual coefficients trong từng block.

\begin{figure}[H]
\centering
\begin{tikzpicture}
\matrix[matrix of nodes,
  nodes={rectangle, draw=LineGray, minimum width=12mm, minimum height=9mm, align=center, font=\small},
  column sep=-\pgflinewidth,
  row sep=-\pgflinewidth] {
0 & 1 & 2 & 3 \\
4 & 5 & 6 & 7 \\
8 & 9 & 10 & 11 \\
12 & 13 & 14 & 15 \\
};
\end{tikzpicture}
\caption{Một macroblock luma có thể được chia thành 16 block 4x4, đánh số 0 đến 15.}
\label{fig:block-index-grid}
\end{figure}

Ví dụ, nếu hệ thống nói vị trí:

\begin{center}
\texttt{(4960, 7, 13)}
\end{center}

thì có thể hiểu là:

\begin{itemize}[leftmargin=*]
  \item macroblock global index 4960,
  \item block 4x4 số 7,
  \item coefficient index 13 trong block đó.
\end{itemize}

\subsection{Một block 4x4 sau khi parse trông như thế nào?}

Giả sử block 7 trong macroblock 4960 có coefficients như sau:

\begin{center}
\begin{tabular}{|c|c|c|c|}
\hline
0 & 5 & -4 & 0 \\
\hline
6 & -1 & 0 & 0 \\
\hline
8 & 0 & -1 & 1 \\
\hline
0 & 0 & 0 & -1 \\
\hline
\end{tabular}
\end{center}

Hệ thống cũng lưu thêm metadata CAVLC cho block này:

\begin{center}
\begin{tabular}{lp{90mm}}
\toprule
Trường nội bộ & Giá trị minh họa \\
\midrule
\texttt{macroblock\_id} & \texttt{4960} \\
\texttt{block\_id} & \texttt{7} \\
\texttt{nC} & \texttt{3}, ngữ cảnh láng giềng dùng khi decode/encode CAVLC. \\
\texttt{bit\_offset} & \texttt{18422}, vị trí bắt đầu block trong RBSP/slice. \\
\texttt{bit\_length} & \texttt{21}, độ dài mã CAVLC ban đầu của block. \\
\texttt{trailing\_ones} & \texttt{2}, số T1 sign flags thật mà block đang có. \\
\texttt{t1\_positions} & \texttt{[13, 15]}, coefficient index có thể là T1 sign positions. \\
\bottomrule
\end{tabular}
\end{center}

Sau bước này, hệ thống chưa nhúng gì cả. Nó mới chỉ biết: ``block này có vài
coefficient có vẻ đáng xem xét''.

\subsection{Bảng candidate positions mẫu}

Từ nhiều macroblock và block, safety filter tạo ra danh sách candidate:

\begin{center}
\begin{tabular}{rrrrlp{38mm}}
\toprule
Candidate & Frame & MB & Block & Coeff & Trạng thái ban đầu \\
\midrule
0 & 12 & 4960 & 7 & 13 & T1 sign candidate. \\
1 & 12 & 4960 & 7 & 15 & T1 sign candidate nhưng cùng block với candidate 0. \\
2 & 13 & 5356 & 2 & 11 & T1 sign candidate. \\
3 & 14 & 5752 & 9 & 14 & T1 sign candidate. \\
\bottomrule
\end{tabular}
\end{center}

Sau đó per-block limit có thể giữ candidate 0 và loại candidate 1, vì cả hai
nằm trong cùng \texttt{(MB, block)}. Đây là lý do số candidate giảm trước khi
thành operating positions.

\subsection{Một bước patchability check mẫu}

Với candidate \texttt{(4960, 7, 13)}, hệ thống thử sửa dấu coefficient rồi
re-encode block bằng CAVLC:

\begin{center}
\begin{tabular}{lll}
\toprule
Kiểm tra & Kết quả minh họa & Ý nghĩa \\
\midrule
Decode block gốc & pass & Parser đọc được block với nC đúng. \\
Encode lại block gốc & 21 bit & Khớp độ dài ban đầu. \\
Sửa coefficient & \(-1 \rightarrow +1\) & Nhúng bit 1. \\
Encode block mới & 21 bit & Độ dài không đổi, có thể patch. \\
Forward round-trip & pass & Decode block mới vẫn ra coefficient mong muốn. \\
Ancestor/retro check & pass & Không làm block trước nuốt nhầm bit của block này. \\
\bottomrule
\end{tabular}
\end{center}

Nếu mọi dòng đều pass, candidate này mới được xem là patchable. Nếu chỉ một
dòng fail, hệ thống loại block hoặc vị trí đó khỏi capacity thật.

\subsection{Từ candidate đến operating position}

Sau nhiều bước lọc, danh sách có thể rút lại như sau:

\begin{center}
\begin{tabular}{lrrrr}
\toprule
Tầng & Số vị trí minh họa & Candidate 0 & Candidate 1 & Candidate 2 \\
\midrule
Raw safe & 433,256 & giữ & giữ & giữ \\
Patchable usable & 1,391 & giữ & loại & giữ \\
Validated pool & 1,316 & giữ & loại & giữ \\
Operating positions & 1,232 & giữ & loại & giữ \\
\bottomrule
\end{tabular}
\end{center}

Khi tới operating positions, mỗi vị trí đã đủ điều kiện để mang một bit payload
trong benchmark-grade path.

\section{Residual Coefficient Là Gì?}

H.264 thường không lưu trực tiếp từng pixel. Nó dự đoán nội dung của một block,
rồi lưu phần sai khác giữa dự đoán và dữ liệu thật. Phần sai khác này gọi là
\textbf{residual}. Sau biến đổi và lượng tử hóa, residual trở thành một dãy số
gọi là \textbf{residual coefficients}.

Ví dụ một block 4x4 có thể có 16 coefficient:

\begin{center}
\begin{tikzpicture}
\matrix[matrix of nodes,
        nodes={coeff},
        column sep=-\pgflinewidth,
        row sep=-\pgflinewidth] (m) {
0 & 5 & -4 & 0 \\
6 & -1 & 0 & 0 \\
8 & 0 & -1 & 1 \\
0 & 0 & 0 & -1 \\
};
\end{tikzpicture}
\end{center}

Hệ thống không sửa tất cả các số này. Nó quan tâm đặc biệt tới các coefficient
\(\pm1\) ở cuối dãy non-zero của CAVLC, vì các coefficient này có thể được mã
hóa như \textbf{trailing-one sign flags}.

\section{CAVLC Và Trailing Ones}

\subsection{CAVLC là gì?}

CAVLC là viết tắt của:

\begin{center}
\textbf{Context-Adaptive Variable Length Coding}
\end{center}

Đây là cách H.264 baseline mã hóa residual coefficients. Điểm quan trọng là
\textbf{variable length}: giá trị khác nhau có thể tạo ra chuỗi bit dài ngắn
khác nhau.

Ví dụ minh họa:

\begin{center}
\begin{tabular}{ccc}
\toprule
Giá trị coefficient & Mã minh họa & Độ dài \\
\midrule
1 & \texttt{1} & 1 bit \\
4 & \texttt{001011} & 6 bit \\
5 & \texttt{0001110} & 7 bit \\
\bottomrule
\end{tabular}
\end{center}

Bảng trên chỉ minh họa ý tưởng, không phải bảng mã H.264 đầy đủ. Ý chính là:
nếu sửa coefficient bừa bãi, độ dài CAVLC có thể đổi. Khi độ dài đổi, việc vá
trở lại bitstream trở nên nguy hiểm.

\subsection{Trailing ones là gì?}

Trong CAVLC, tối đa ba coefficient cuối có trị tuyệt đối bằng 1 có thể được mã
hóa đặc biệt. Chúng được gọi là \textbf{trailing ones}, hay viết tắt là
\textbf{T1}. Với T1, phần dấu \((+/-)\) được mã hóa bằng các bit dấu riêng.

Điều này rất hợp với nhúng tin:

\begin{center}
\begin{tabular}{ccc}
\toprule
Bit cần nhúng & Hành động trực quan & Coefficient sau cùng \\
\midrule
0 & đặt dấu âm & \(-1\) \\
1 & đặt dấu dương & \(+1\) \\
\bottomrule
\end{tabular}
\end{center}

Nếu coefficient thật sự là T1 sign flag, đổi dấu thường là thay đổi rất nhỏ và
có thể giữ nguyên độ dài mã. Vì vậy hệ thống tập trung vào \textbf{T1 sign
embedding}.

\section{Ví Dụ Cụ Thể: Nhúng 4 Bit Vào Một Số Vị Trí}

Giả sử payload có 4 bit đầu:

\begin{center}
\texttt{1 0 1 1}
\end{center}

Hệ thống đã chọn được 4 vị trí T1 an toàn:

\begin{center}
\begin{tabular}{clcc}
\toprule
Bit & Vị trí & Coefficient trước & Coefficient sau \\
\midrule
1 & \texttt{(12, 496, 1)} & \(-1\) & \(+1\) \\
0 & \texttt{(13, 892, 1)} & \(+1\) & \(-1\) \\
1 & \texttt{(14, 1288, 1)} & \(+1\) & \(+1\) \\
1 & \texttt{(15, 1684, 1)} & \(-1\) & \(+1\) \\
\bottomrule
\end{tabular}
\end{center}

Điều quan trọng là mỗi dòng trong bảng không chỉ là ``sửa một số''. Mỗi sửa đổi
phải thỏa mãn:

\begin{itemize}[leftmargin=*]
  \item coefficient đúng là vị trí T1 sign flag;
  \item block sau khi sửa vẫn mã hóa được bằng CAVLC;
  \item độ dài bit của block không phá cấu trúc NAL/slice;
  \item số sửa đổi trong một block không vượt giới hạn;
  \item video sau tái tạo vẫn decode được;
  \item bit có thể đọc lại từ stego video.
\end{itemize}

\section{Pipeline Chọn Vị Trí Nhúng}

Đây là phần quan trọng nhất của hệ thống. Một vị trí chỉ được dùng nếu đi qua
nhiều cửa kiểm tra.

\begin{figure}[H]
\centering
\resizebox{\textwidth}{!}{%
\begin{tikzpicture}[node distance=8mm and 8mm]
  \node[widebox] (parse) {Parse H.264\\lấy coefficients, offsets, nC};
  \node[widebox, right=of parse] (t1) {Tìm T1 thật\\không nhầm với \(\pm1\) thường};
  \node[widebox, right=of t1] (safe) {Safety filter\\raw safe positions};
  \node[widebox, below=of safe] (patch) {Patchability pruning\\block phải vá được};
  \node[widebox, left=of patch] (limit) {Per-block limit\\thường 1 bit/block};
  \node[widebox, left=of limit] (validate) {Validation/quality\\PSNR, FFmpeg, SEC1 pool};
  \node[good, below=of limit] (op) {Operating positions\\vị trí dùng thật};

  \draw[arrow] (parse) -- (t1);
  \draw[arrow] (t1) -- (safe);
  \draw[arrow] (safe) -- (patch);
  \draw[arrow] (patch) -- (limit);
  \draw[arrow] (limit) -- (validate);
  \draw[arrow] (validate) -- (op);
\end{tikzpicture}
}
\caption{Pipeline chọn vị trí nhúng từ H.264 bitstream đến operating positions.}
\label{fig:position-pipeline}
\end{figure}

\subsection{Safety filter}

Safety filter là lớp lọc đầu tiên. Nó trả lời câu hỏi:

\begin{center}
\textit{``Vị trí này có vẻ an toàn về mặt cấu trúc CAVLC không?''}
\end{center}

Nó kiểm tra các yếu tố như:

\begin{itemize}[leftmargin=*]
  \item vị trí có phải trailing-one sign flag thật không;
  \item sửa dấu có giữ được kiểu mã hóa mong muốn không;
  \item block không nằm trong nhóm đã biết là không patchable;
  \item thay đổi không làm zero/non-zero pattern nguy hiểm;
  \item magnitude và số sửa đổi không vượt ngưỡng.
\end{itemize}

\subsection{Patchability-aware embedding}

Patchability-aware embedding nghĩa là:

\begin{center}
\textit{``Chỉ nhúng vào những vị trí mà sau khi sửa, block vẫn có thể được vá
lại vào bitstream H.264 hợp lệ.''}
\end{center}

Đây là một khái niệm rất quan trọng. Trong H.264/CAVLC, sửa một coefficient có
thể làm block sau khi mã hóa dài hơn hoặc ngắn hơn. Nếu độ dài thay đổi, phần
bit phía sau trong slice có thể bị lệch. Khi đó video có thể hỏng, parser có
thể đọc sai, hoặc verifier không đọc lại được payload.

\begin{figure}[H]
\centering
\resizebox{\textwidth}{!}{%
\begin{tikzpicture}[node distance=7mm and 10mm]
  \node[widebox] (orig1) {Block gốc\\21 bit CAVLC};
  \node[widebox, right=of orig1] (mod1) {Block sau sửa\\23 bit CAVLC};
  \node[bad, right=of mod1] (bad) {Không patchable\\lệch bitstream};

  \node[widebox, below=15mm of orig1] (orig2) {Block gốc\\21 bit CAVLC};
  \node[widebox, right=of orig2] (mod2) {Block sau sửa\\21 bit CAVLC};
  \node[good, right=of mod2] (good) {Patchable\\có thể vá lại};

  \draw[arrow] (orig1) -- node[smallnote, above]{sửa coefficient} (mod1);
  \draw[arrow] (mod1) -- (bad);
  \draw[arrow] (orig2) -- node[smallnote, above]{sửa T1 sign} (mod2);
  \draw[arrow] (mod2) -- (good);
\end{tikzpicture}
}
\caption{Patchability: sửa xong phải vá được vào đúng vị trí trong bitstream.}
\label{fig:patchability}
\end{figure}

Ảnh hưởng của patchability:

\begin{itemize}[leftmargin=*]
  \item \textbf{Tăng độ đúng}: video sau cùng ít bị hỏng bitstream.
  \item \textbf{Tăng khả năng đọc lại}: bit đã nhúng có khả năng sống qua reconstruction.
  \item \textbf{Giảm capacity}: nhiều vị trí nhìn có vẻ dùng được sẽ bị loại.
  \item \textbf{Làm claim paper chặt hơn}: capacity được báo gần với khả năng dùng thật.
\end{itemize}

\subsection{Per-block limit}

Hệ thống thường giới hạn số bit nhúng trong mỗi block, ví dụ:

\begin{center}
\texttt{max\_modifications\_per\_block = 1}
\end{center}

Nghĩa là một block 4x4 dù có nhiều vị trí có vẻ dùng được, hệ thống cũng chỉ
chọn tối đa một vị trí. Điều này giúp giảm méo cục bộ và làm reconstruction ổn
định hơn. Đổi lại, capacity giảm.

\subsection{Quality guard}

Quality guard đảm bảo video sau nhúng không bị xấu quá mức. Hệ thống dùng các
metric như PSNR/SSIM, trong đó ngưỡng quan trọng của benchmark hiện tại là:

\begin{center}
\texttt{min\_modified\_frame\_psnr >= 40 dB}
\end{center}

Nếu một cấu hình nhúng làm frame tụt dưới ngưỡng, cấu hình đó không được coi là
operating point mạnh.

\section{Capacity: Vì Sao Không Thể Chỉ Đếm Candidate?}

Một điểm mạnh của hệ thống là nó không nói ``video có bao nhiêu coefficient thì
có bấy nhiêu bit''. Nó chia capacity thành nhiều lớp.

\begin{figure}[H]
\centering
\begin{tikzpicture}[node distance=8mm]
  \node[widebox, text width=78mm] (raw) {Raw safe bits\\vị trí qua safety filter};
  \node[widebox, text width=68mm, below=of raw] (patch) {Patchable usable bits\\vị trí vá được và bị giới hạn theo block};
  \node[widebox, text width=58mm, below=of patch] (valid) {Validated pool bits\\vị trí qua validation/quality};
  \node[good, text width=48mm, below=of valid] (op) {Operating bits\\vị trí dùng thật trong benchmark};

  \draw[arrow] (raw) -- node[smallnote, right]{loại block không vá được} (patch);
  \draw[arrow] (patch) -- node[smallnote, right]{kiểm tra thêm} (valid);
  \draw[arrow] (valid) -- node[smallnote, right]{khóa operating point} (op);
\end{tikzpicture}
\caption{Capacity funnel: càng xuống dưới càng gần capacity thật.}
\label{fig:capacity-funnel}
\end{figure}

Ví dụ minh họa:

\begin{center}
\begin{tabular}{lrp{70mm}}
\toprule
Lớp capacity & Số lượng ví dụ & Ý nghĩa \\
\midrule
Raw safe bits & 433,256 & Vị trí có vẻ an toàn về cấu trúc CAVLC. \\
Patchable usable bits & 1,391 & Vị trí còn lại sau khi xét patchability và giới hạn theo block. \\
Validated pool bits & 1,316 & Vị trí đã qua validation/quality. \\
Operating bits & 1,232 & Vị trí thật sự dùng ở operating point. \\
\bottomrule
\end{tabular}
\end{center}

Điểm quan trọng là: \textbf{raw safe bits không phải capacity thật}. Capacity
thật hơn nằm ở patchable/validated/operating layers.

\section{Payload Được Chuẩn Bị Như Thế Nào?}

Về mặt video, hệ thống chỉ cần một chuỗi bit để nhúng. Chuỗi bit này đến từ
payload blob:

\begin{center}
\texttt{[4B message length][message][129B Groth16 proof]}
\end{center}

Ví dụ message benchmark:

\begin{center}
\texttt{b"ZK-bench-v1.0!"}
\end{center}

Nếu message dài 13 byte và proof dài 129 byte, payload chưa chaos là:

\begin{center}
\texttt{4 + 13 + 129 = 146 bytes = 1168 bits}
\end{center}

Trong một số benchmark có chaos transform, payload operating có thể mở rộng lên
khoảng 1232 bit.

\section{Payload Được Nhúng Vào Video Như Thế Nào?}

Sau khi có payload bits và operating positions, hệ thống ghép từng bit với từng
vị trí:

\begin{figure}[H]
\centering
\resizebox{\textwidth}{!}{%
\begin{tikzpicture}[node distance=8mm and 9mm]
  \node[crypto] (payload) {Payload blob\\\texttt{[len][msg][proof]}};
  \node[widebox, right=of payload] (bits) {Đổi thành bit\\\(b_0,b_1,\ldots,b_n\)};
  \node[widebox, below=of bits] (positions) {Operating positions\\\(p_0,p_1,\ldots,p_n\)};
  \node[highlight, right=of bits] (map) {Ánh xạ\\bit \(b_i\) vào vị trí \(p_i\)};
  \node[widebox, right=of map] (modify) {Sửa T1 sign\\coefficient \(\pm1\)};
  \node[good, below=of modify] (modmap) {Modified coefficient map};

  \draw[arrow] (payload) -- (bits);
  \draw[arrow] (bits) -- (map);
  \draw[arrow] (positions) -- (map);
  \draw[arrow] (map) -- (modify);
  \draw[arrow] (modify) -- (modmap);
\end{tikzpicture}
}
\caption{Mỗi payload bit được đưa vào một vị trí coefficient đã chọn.}
\label{fig:embed-bits}
\end{figure}

\subsection{Ví dụ ánh xạ payload bits vào positions}

Giả sử payload sau khi pack được đổi thành các bit đầu tiên:

\begin{center}
\texttt{1 0 1 1 0 0 1 0}
\end{center}

Và operating positions đầu tiên là:

\begin{center}
\begin{tabular}{rrrrcc}
\toprule
Bit index & Payload bit & Frame & MB & Block & Coeff \\
\midrule
0 & 1 & 12 & 4960 & 7 & 13 \\
1 & 0 & 13 & 5356 & 2 & 11 \\
2 & 1 & 14 & 5752 & 9 & 14 \\
3 & 1 & 15 & 6148 & 1 & 15 \\
4 & 0 & 16 & 6544 & 6 & 12 \\
5 & 0 & 17 & 6940 & 5 & 13 \\
6 & 1 & 18 & 7336 & 8 & 15 \\
7 & 0 & 19 & 7732 & 3 & 11 \\
\bottomrule
\end{tabular}
\end{center}

Sau đó hệ thống đi từng dòng:

\begin{center}
\begin{tabular}{rrrrccc}
\toprule
Bit & Frame & MB & Block & Coeff & Trước & Sau \\
\midrule
1 & 12 & 4960 & 7 & 13 & \(-1\) & \(+1\) \\
0 & 13 & 5356 & 2 & 11 & \(+1\) & \(-1\) \\
1 & 14 & 5752 & 9 & 14 & \(+1\) & \(+1\) \\
1 & 15 & 6148 & 1 & 15 & \(-1\) & \(+1\) \\
0 & 16 & 6544 & 6 & 12 & \(-1\) & \(-1\) \\
0 & 17 & 6940 & 5 & 13 & \(+1\) & \(-1\) \\
1 & 18 & 7336 & 8 & 15 & \(-1\) & \(+1\) \\
0 & 19 & 7732 & 3 & 11 & \(+1\) & \(-1\) \\
\bottomrule
\end{tabular}
\end{center}

Bảng này là hình ảnh trực quan nhất của quá trình nhúng: payload bit được biến
thành dấu của một coefficient T1 ở một block cụ thể.

Ở mức một block, có thể hình dung như sau:

\begin{figure}[H]
\centering
\resizebox{\textwidth}{!}{%
\begin{tabular}{ccc}
\toprule
\textbf{Block trước nhúng} &
\textbf{Hành động} &
\textbf{Block sau nhúng} \\
\midrule
\begin{tabular}{|c|c|c|c|}
\hline
0 & 5 & -4 & 0 \\
\hline
6 & -1 & 0 & 0 \\
\hline
8 & 0 & -1 & 1 \\
\hline
0 & 0 & 0 & \cellcolor{SoftYellow}\(-1\) \\
\hline
\end{tabular}
&
\begin{tabular}{c}
Payload bit \(=1\) \\
chọn T1 sign position \\
đặt coefficient thành \(+1\)
\end{tabular}
&
\begin{tabular}{|c|c|c|c|}
\hline
0 & 5 & -4 & 0 \\
\hline
6 & -1 & 0 & 0 \\
\hline
8 & 0 & -1 & 1 \\
\hline
0 & 0 & 0 & \cellcolor{CoreGreen!12}\(+1\) \\
\hline
\end{tabular}
\\
\bottomrule
\end{tabular}
}
\caption{Ví dụ trực quan: nhúng một bit bằng cách đổi dấu một coefficient T1.}
\label{fig:block-example}
\end{figure}

\section{Reconstruction: Làm Sao Video Sau Nhúng Vẫn Là H.264 Hợp Lệ?}

Sau khi sửa coefficient trong bộ nhớ, hệ thống phải tạo lại file stego H.264.
Đây là bước reconstruction.

\begin{figure}[H]
\centering
\resizebox{\textwidth}{!}{%
\begin{tikzpicture}[node distance=8mm and 8mm]
  \node[widebox] (modcoeff) {Modified coefficients\\trong bộ nhớ};
  \node[widebox, right=of modcoeff] (context) {Reconstruction context\\offsets, nC, T1 override};
  \node[widebox, right=of context] (encode) {Re-encode/patch\\CAVLC blocks/slices};
  \node[widebox, below=of encode] (out) {Stego H.264\\file đầu ra};
  \node[widebox, left=of out] (stats) {Applied block keys\\block thật sự được vá};
  \node[good, left=of stats] (ok) {Post-check\\đủ bit thì pass};

  \draw[arrow] (modcoeff) -- (context);
  \draw[arrow] (context) -- (encode);
  \draw[arrow] (encode) -- (out);
  \draw[arrow] (out) -- (stats);
  \draw[arrow] (stats) -- (ok);
\end{tikzpicture}
}
\caption{Reconstruction biến coefficient đã sửa thành stego H.264 hợp lệ.}
\label{fig:reconstruction}
\end{figure}

Sau reconstruction, hệ thống không tin mù rằng mọi bit đều đã được áp dụng.
Nó kiểm tra lại:

\begin{itemize}[leftmargin=*]
  \item \texttt{requested\_position\_bits}: số vị trí dự định dùng;
  \item \texttt{applied\_position\_bits}: số vị trí thật sự được vá vào bitstream;
  \item \texttt{bits\_required}: số bit payload cần nhúng.
\end{itemize}

Nếu \texttt{applied\_position\_bits < bits\_required}, hệ thống báo fail ở
\texttt{post\_reconstruct\_application}. Đây là một điểm rất quan trọng: hệ
thống đo capacity sau khi bitstream đã được tái tạo, không chỉ trước khi tái
tạo.

\subsection{Ví dụ artifact sau khi nhúng}

Sau một lượt nhúng thành công, thư mục output thường có các file dạng:

\begin{center}
\begin{tabular}{lp{88mm}}
\toprule
File & Ý nghĩa \\
\midrule
\texttt{stego.h264} & Video H.264 đã chứa payload. \\
\texttt{stego.h264.positions.json} & Danh sách positions đã dùng để nhúng bit. \\
\texttt{stego.h264.meta.json} & Số bit required/embedded, capacity diagnostics, chaos flag. \\
\texttt{stego.h264.manifest.json} & Manifest có schema version, payload metadata, video hash, proof metadata. \\
\bottomrule
\end{tabular}
\end{center}

Một \texttt{positions.json} rút gọn có thể nhìn như sau:

\begin{center}
\begin{tabular}{rrrr}
\toprule
Index & Macroblock & Block & Coefficient \\
\midrule
0 & 4960 & 7 & 13 \\
1 & 5356 & 2 & 11 \\
2 & 5752 & 9 & 14 \\
3 & 6148 & 1 & 15 \\
\bottomrule
\end{tabular}
\end{center}

Một \texttt{meta.json} rút gọn có thể chứa:

\begin{center}
\begin{tabular}{lr}
\toprule
Trường & Giá trị minh họa \\
\midrule
\texttt{bits\_required} & 1232 \\
\texttt{bits\_embedded} & 1232 \\
\texttt{raw\_safe\_bits} & 433256 \\
\texttt{patchable\_usable\_bits} & 1391 \\
\texttt{requested\_position\_bits} & 1232 \\
\texttt{applied\_position\_bits} & 1232 \\
\bottomrule
\end{tabular}
\end{center}

Nếu \texttt{applied\_position\_bits} nhỏ hơn \texttt{bits\_required}, điều đó có
nghĩa là một phần positions đã không sống qua reconstruction, và output đó
không được xem là operating-point success.

\section{Trích Xuất Payload Diễn Ra Như Thế Nào?}

Trích xuất là pipeline ngược của nhúng.

\begin{figure}[H]
\centering
\resizebox{\textwidth}{!}{%
\begin{tikzpicture}[node distance=8mm and 8mm]
  \node[widebox] (stego) {Stego H.264};
  \node[widebox, right=of stego] (parse) {Parse lại bitstream\\NAL, IDR, block};
  \node[widebox, right=of parse] (positions) {Lấy positions\\từ sidecar hoặc locked set};
  \node[widebox, below=of positions] (read) {Đọc T1 signs\\khôi phục bits};
  \node[widebox, left=of read] (blob) {Ghép lại payload blob};
  \node[crypto, left=of blob] (unpack) {Unpack\\length, message, proof};
  \node[crypto, below=of unpack] (verify) {Verify proof\\Groth16};

  \draw[arrow] (stego) -- (parse);
  \draw[arrow] (parse) -- (positions);
  \draw[arrow] (positions) -- (read);
  \draw[arrow] (read) -- (blob);
  \draw[arrow] (blob) -- (unpack);
  \draw[arrow] (unpack) -- (verify);
\end{tikzpicture}
}
\caption{Pipeline trích xuất: đọc dấu T1 ở đúng positions để dựng lại payload.}
\label{fig:extract}
\end{figure}

\subsection{Ví dụ đọc ngược từ stego video}

Giả sử verifier có \texttt{positions.json} với các dòng đầu:

\begin{center}
\begin{tabular}{rrrr}
\toprule
Index & Macroblock & Block & Coefficient \\
\midrule
0 & 4960 & 7 & 13 \\
1 & 5356 & 2 & 11 \\
2 & 5752 & 9 & 14 \\
3 & 6148 & 1 & 15 \\
\bottomrule
\end{tabular}
\end{center}

Verifier parse \texttt{stego.h264}, đi tới đúng các vị trí đó và đọc dấu của
coefficient:

\begin{center}
\begin{tabular}{rrrrcc}
\toprule
Index & MB & Block & Coeff & Coefficient đọc được & Bit khôi phục \\
\midrule
0 & 4960 & 7 & 13 & \(+1\) & 1 \\
1 & 5356 & 2 & 11 & \(-1\) & 0 \\
2 & 5752 & 9 & 14 & \(+1\) & 1 \\
3 & 6148 & 1 & 15 & \(+1\) & 1 \\
\bottomrule
\end{tabular}
\end{center}

Các bit này được ghép lại thành payload blob:

\begin{center}
\texttt{1011...}
\end{center}

Sau khi đủ số bit, hệ thống đọc 4 byte đầu để biết độ dài message, lấy phần
message, lấy phần proof, rồi gọi verifier ZKP. Nếu một vài bit đầu header sai,
toàn bộ quá trình unpack có thể hỏng vì độ dài message bị đọc sai. Đây là lý
do các thử nghiệm blind-core sau này phải quan tâm rất nhiều tới header
reliability.

Với verifier hiện tại:

\begin{itemize}[leftmargin=*]
  \item \textbf{Strict non-blind}: cần cover/original hoặc thông tin vị trí tương đương.
  \item \textbf{Sidecar-assisted near-blind}: không cần cover gốc nhưng cần sidecar như manifest/positions.
  \item \textbf{Blind-core}: hướng nghiên cứu, chưa phải claim chính của baseline hiện tại.
\end{itemize}

\section{Cơ Chế ZKP: Proof Được Sinh Từ Đâu Và Sinh Như Thế Nào?}

ZKP không chọn vị trí nhúng trong video. ZKP tạo ra một proof mật mã để chứng
minh rằng message được nhúng có liên hệ đúng với một secret key, nhưng không
tiết lộ secret key đó. Sau đó proof được nén thành bytes và đưa cho video
embedding core giấu vào H.264 như dữ liệu bình thường.

\subsection{Mệnh đề mà proof chứng minh}

Circuit chính nằm trong:

\begin{center}
\texttt{circuits/payload\_verify.circom}
\end{center}

Circuit tên là \texttt{PayloadVerify}. Nó chứng minh quan hệ:

\[
  \texttt{commitment}
  =
  \mathrm{SHA256}(\texttt{payload\_hash} \parallel \texttt{secret})
\]

Trong đó:

\[
  \texttt{payload\_hash}
  =
  \mathrm{SHA256}(\texttt{message})
\]

Nói bằng lời:

\begin{center}
\textit{Người tạo proof biết một secret key sao cho khi ghép secret key đó với
hash của message, rồi băm SHA256, ta nhận được commitment công khai.}
\end{center}

Điểm quan trọng là secret key không được nhúng vào video và không xuất hiện
trong public signals. Secret key chỉ đi vào witness khi tạo proof.

\subsection{Public input và private input}

Circuit có các input sau:

\begin{center}
\begin{tabular}{lll}
\toprule
Loại input & Tên & Ý nghĩa \\
\midrule
Public & \texttt{payload\_hash[256]} & 256 bit của \(\mathrm{SHA256(message)}\). \\
Public & \texttt{commitment[256]} & 256 bit commitment mong đợi. \\
Public & \texttt{payload\_length} & Độ dài message theo byte. \\
Private & \texttt{secret[256]} & 256 bit secret key, tức 32 byte. \\
\bottomrule
\end{tabular}
\end{center}

Public input là phần verifier được phép biết. Private input là phần prover biết
nhưng không lộ ra. Đây là phần ``zero-knowledge'' của hệ thống.

\subsection{Ví dụ dữ liệu trước khi tạo proof}

Giả sử message là:

\begin{center}
\texttt{b"ZK-bench-v1.0!"}
\end{center}

và secret key là 32 byte:

\begin{center}
\texttt{00 01 02 03 ... 1f}
\end{center}

Python bridge trong \texttt{src/zk\_proof.py} tạo dữ liệu circuit như sau:

\begin{center}
\begin{tabular}{lp{95mm}}
\toprule
Bước & Giá trị tạo ra \\
\midrule
1 & \texttt{payload\_hash = SHA256(message)} \\
2 & \texttt{commitment = SHA256(payload\_hash || secret\_key)} \\
3 & Chuyển \texttt{payload\_hash} thành 256 bit. \\
4 & Chuyển \texttt{commitment} thành 256 bit. \\
5 & Chuyển \texttt{secret\_key} thành 256 bit private input. \\
6 & Ghi thêm \texttt{payload\_length = len(message)}. \\
\bottomrule
\end{tabular}
\end{center}

Trong code, logic này tương ứng với hàm
\texttt{\_build\_circuit\_input(payload\_bytes, secret\_key)}.

\subsection{Workflow sinh proof}

Hệ thống dùng \texttt{snarkjs} và Groth16. Các artifact quan trọng nằm trong
\texttt{circuits/build/}:

\begin{center}
\begin{tabular}{lp{88mm}}
\toprule
Artifact & Vai trò \\
\midrule
\texttt{payload\_verify.wasm} & Chạy circuit để tạo witness từ input. \\
\texttt{generate\_witness.js} & Script Node.js sinh witness. \\
\texttt{proving\_key.zkey} & Proving key dùng để tạo Groth16 proof. \\
\texttt{verification\_key.json} & Verification key dùng để verify proof. \\
\bottomrule
\end{tabular}
\end{center}

Pipeline tạo proof:

\begin{figure}[H]
\centering
\resizebox{\textwidth}{!}{%
\begin{tikzpicture}[node distance=8mm and 8mm]
  \node[crypto] (message) {Message bytes};
  \node[crypto, right=of message] (secret) {Secret key\\32 bytes};
  \node[crypto, below=of message] (hash) {\(\mathrm{SHA256(message)}\)\\payload hash};
  \node[crypto, right=of hash] (commit) {\(\mathrm{SHA256(hash || secret)}\)\\commitment};
  \node[crypto, below=of hash] (input) {Circuit input JSON\\public + private};
  \node[crypto, right=of input] (witness) {Witness\\\texttt{witness.wtns}};
  \node[crypto, right=of witness] (proof) {Groth16 proof\\snarkjs JSON};
  \node[crypto, below=of proof] (bytes) {Compressed proof\\129 bytes};
  \node[widebox, left=of bytes] (pack) {Payload blob\\\texttt{[4B len][message][proof]}};
  \node[widebox, left=of pack] (video) {Embed blob bits\\into H.264};

  \draw[arrow] (message) -- (hash);
  \draw[arrow] (message) -- (input);
  \draw[arrow] (secret) -- (commit);
  \draw[arrow] (hash) -- (commit);
  \draw[arrow] (hash) -- (input);
  \draw[arrow] (commit) -- (input);
  \draw[arrow] (secret) -- (input);
  \draw[arrow] (input) -- node[smallnote, above]{Node.js WASM} (witness);
  \draw[arrow] (witness) -- node[smallnote, above]{snarkjs groth16 prove} (proof);
  \draw[arrow] (proof) -- (bytes);
  \draw[arrow] (bytes) -- (pack);
  \draw[arrow] (pack) -- (video);
\end{tikzpicture}
}
\caption{Workflow tạo proof: message và secret tạo witness, witness tạo Groth16 proof, proof được nén rồi nhúng vào video.}
\label{fig:zkp-proof-generation}
\end{figure}

\subsection{Witness là gì?}

Witness là tập giá trị làm cho circuit thỏa mãn. Trong hệ thống này, witness
bao gồm secret key và các giá trị trung gian khi circuit tính:

\[
  \mathrm{SHA256}(\texttt{payload\_hash} \parallel \texttt{secret})
\]

Nếu secret không đúng, circuit không thể tạo witness hợp lệ cho commitment đã
công khai. Nếu witness không hợp lệ, Groth16 proof sẽ không verify.

Trong code, witness được tạo bằng:

\begin{center}
\texttt{node generate\_witness.js payload\_verify.wasm input.json witness.wtns}
\end{center}

Sau khi tạo proof, các file tạm chứa secret/witness/proof/public được xóa.
Điểm này quan trọng vì input JSON có chứa secret key ở dạng bit.

\subsection{Proof JSON và proof bytes khác nhau thế nào?}

\texttt{snarkjs} tạo proof ở dạng JSON, thường có các phần:

\begin{center}
\begin{tabular}{lp{85mm}}
\toprule
Trường proof & Ý nghĩa \\
\midrule
\texttt{pi\_a} & Điểm elliptic curve trong group G1. \\
\texttt{pi\_b} & Điểm elliptic curve trong group G2. \\
\texttt{pi\_c} & Điểm elliptic curve trong group G1. \\
\texttt{protocol} & \texttt{groth16}. \\
\texttt{curve} & \texttt{bn128}. \\
\bottomrule
\end{tabular}
\end{center}

JSON này quá dài để nhúng trực tiếp. Vì vậy hệ thống nén proof thành 129 byte:

\begin{center}
\begin{tabular}{lr}
\toprule
Thành phần & Kích thước \\
\midrule
\texttt{pi\_a.x} & 32 byte \\
\texttt{pi\_b.x[0]} & 32 byte \\
\texttt{pi\_b.x[1]} & 32 byte \\
\texttt{pi\_c.x} & 32 byte \\
sign flags cho các tọa độ y & 1 byte \\
\midrule
Tổng & 129 byte \\
\bottomrule
\end{tabular}
\end{center}

Hàm tương ứng là:

\begin{center}
\texttt{proof\_to\_bytes(proof\_dict)}
\end{center}

Khi verify, hệ thống gọi:

\begin{center}
\texttt{bytes\_to\_proof(proof\_bytes)}
\end{center}

để dựng lại proof JSON từ 129 byte này.

\subsection{Payload blob cuối cùng được nhúng vào video}

Sau khi có proof bytes, hệ thống pack message và proof thành một blob:

\begin{center}
\texttt{[4 bytes message length][message bytes][129 bytes proof]}
\end{center}

Ví dụ:

\begin{center}
\begin{tabular}{lrr}
\toprule
Thành phần & Kích thước & Ghi chú \\
\midrule
Message length header & 4 byte & Big-endian integer. \\
Message & 13 byte & Ví dụ \texttt{ZK-bench-v1.0!}. \\
Compressed proof & 129 byte & Groth16 BN128 proof. \\
\midrule
Tổng & 146 byte & 1168 bit trước chaos. \\
\bottomrule
\end{tabular}
\end{center}

Đây là dữ liệu mà video embedding core thật sự nhìn thấy. Từ góc nhìn của phần
nhúng video, ZKP proof chỉ là 129 byte cần được giấu và đọc lại chính xác.

\subsection{Verify sau khi extract diễn ra thế nào?}

Sau khi verifier đọc lại payload blob từ stego video, hệ thống làm ngược lại:

\begin{enumerate}[leftmargin=*]
  \item Đọc 4 byte đầu để biết độ dài message.
  \item Tách message bytes và 129 proof bytes.
  \item Dựng lại proof JSON bằng \texttt{bytes\_to\_proof()}.
  \item Tính lại \texttt{payload\_hash = SHA256(message)}.
  \item Tính lại \texttt{commitment = SHA256(payload\_hash || secret\_key)}.
  \item Tạo public signals gồm \texttt{payload\_hash}, \texttt{commitment}, \texttt{payload\_length}.
  \item Gọi \texttt{snarkjs groth16 verify} với \texttt{verification\_key.json}.
\end{enumerate}

\begin{figure}[H]
\centering
\resizebox{\textwidth}{!}{%
\begin{tikzpicture}[node distance=8mm and 8mm]
  \node[widebox] (stego) {Extracted blob\\from stego H.264};
  \node[crypto, right=of stego] (unpack) {Unpack blob\\message + proof bytes};
  \node[crypto, right=of unpack] (reproof) {Rebuild proof JSON\\from 129 bytes};
  \node[crypto, below=of unpack] (public) {Rebuild public signals\\hash, commitment, length};
  \node[crypto, right=of public] (verify) {snarkjs\\groth16 verify};
  \node[good, right=of verify] (ok) {Valid proof};

  \draw[arrow] (stego) -- (unpack);
  \draw[arrow] (unpack) -- (reproof);
  \draw[arrow] (unpack) -- (public);
  \draw[arrow] (reproof) -- (verify);
  \draw[arrow] (public) -- (verify);
  \draw[arrow] (verify) -- (ok);
\end{tikzpicture}
}
\caption{Workflow verify: sau khi extract đúng bytes, hệ thống dựng lại proof và public signals rồi gọi Groth16 verifier.}
\label{fig:zkp-verify}
\end{figure}

\subsection{ZKP và video embedding chia trách nhiệm thế nào?}

Hai phần này có vai trò khác nhau:

\begin{center}
\begin{tabular}{p{0.32\linewidth}p{0.58\linewidth}}
\toprule
Thành phần & Trách nhiệm \\
\midrule
ZKP layer & Chứng minh message có commitment hợp lệ với secret key mà không lộ secret key. \\
Video embedding core & Giấu chính xác payload bytes vào H.264 và đọc lại đúng payload bytes đó. \\
Verifier & Kiểm tra rằng bytes đọc lại từ video tạo thành proof hợp lệ. \\
\bottomrule
\end{tabular}
\end{center}

Nói cách khác:

\begin{center}
\textbf{ZKP trả lời: payload có quan hệ mật mã hợp lệ không?}\\
\textbf{Video embedding trả lời: payload có được giấu và đọc lại đúng không?}
\end{center}

\section{Workflow Thực Tế Khi Gọi \texttt{embed()}}

Khi gọi public API \texttt{embed()}, pipeline thực tế đi qua các bước sau:

\begin{enumerate}[leftmargin=*]
  \item Kiểm tra input: video tồn tại, \texttt{circuits\_dir} tồn tại, message không rỗng, secret key dài 32 byte.
  \item Phân tích stream profile: all-intra hay inter-coded. Baseline mạnh nhất hiện tại là all-intra H.264/CAVLC.
  \item Tạo Groth16 proof cho message và secret key.
  \item Pack payload thành \texttt{[4B length][message][proof]}.
  \item Nếu bật chaos, scramble payload bits và shuffle thứ tự positions.
  \item Parse video H.264 để lấy coefficients, CAVLC offsets, nC map, NAL length map, T1 override map.
  \item Chạy CAVLC safety filter để lấy raw safe positions.
  \item Nếu dùng locked operating point, nạp precomputed positions đã được validate.
  \item Nếu không dùng locked positions, chạy FFmpeg validation tùy chọn, patchability pruning, và per-block limit.
  \item Nhúng payload bits vào positions bằng T1 sign modification.
  \item Reconstruct stego H.264.
  \item Kiểm tra applied positions sau reconstruction.
  \item Ghi \texttt{positions.json}, \texttt{meta.json}, và \texttt{manifest.json}.
\end{enumerate}

\section{Tại Sao Hệ Thống Hiện Mạnh Nhất Ở All-Intra H.264/CAVLC?}

All-intra nghĩa là các frame chủ yếu được mã hóa độc lập, giống nhiều I-frame.
Điều này hợp với hệ thống vì:

\begin{itemize}[leftmargin=*]
  \item sửa một frame ít gây cascade sang frame khác;
  \item residual blocks trong IDR/I-frame dễ kiểm soát hơn;
  \item CAVLC baseline có cấu trúc trailing-one phù hợp với sign-bit embedding;
  \item benchmark có thể khóa operating point ổn định hơn.
\end{itemize}

Với GOP lớn hoặc inter-coded video, P-frame/B-frame phụ thuộc dự đoán từ frame
khác. Một thay đổi nhỏ có thể ảnh hưởng rộng hơn. Vì vậy inter-coded support là
hướng cần phát triển thêm, không phải path mạnh nhất hiện tại.

\section{Tóm Tắt Vai Trò Của Các Thành Phần}

\begin{center}
\begin{tabular}{p{0.28\linewidth}p{0.62\linewidth}}
\toprule
Thành phần & Vai trò \\
\midrule
\texttt{src/embedder.py} &
Public API. Tạo proof, pack payload, nạp analysis, chọn positions, nhúng, reconstruct, ghi sidecar. \\

\texttt{src/core/stego.py} &
Safety filter và payload embed/extract logic. Đây là nơi tìm T1 positions và sửa coefficient. \\

\texttt{src/bitstream/cavlc.py} &
CAVLC encoder/decoder. Dùng để hiểu và tái mã hóa residual blocks. \\

\texttt{src/bitstream/h264.py} &
H.264 parser. Dùng để bóc NAL, slice, macroblock, block, coefficient. \\

\texttt{src/bitstream/bitstream\_ops.py} &
Patchability validation và reconstruction. Đây là phần đảm bảo thay đổi có thể sống trong bitstream thật. \\

\texttt{src/zk\_proof.py} &
Tạo, serialize, pack, unpack, và verify Groth16 proof. \\

\texttt{benchmark/\_common.py} &
Đo capacity layers và tái dùng analysis cache cho benchmark. \\
\bottomrule
\end{tabular}
\end{center}

\section{Điểm Cần Nhớ}

Điểm tinh túy của hệ thống không phải là ``đổi coefficient''. Điểm tinh túy là:

\begin{center}
\textit{Chọn đúng coefficient mà sau khi đổi, H.264 bitstream vẫn vá được, video
vẫn decode được, chất lượng vẫn đạt, và payload vẫn đọc lại được.}
\end{center}

Vì thế, các thuật ngữ quan trọng nhất là:

\begin{itemize}[leftmargin=*]
  \item \textbf{CAVLC-safe}: vị trí không phá cấu trúc mã hóa coefficient.
  \item \textbf{Patchable}: block sau sửa có thể vá lại vào bitstream.
  \item \textbf{Reconstruct-aware}: sau reconstruction vẫn kiểm tra số bit sống thật.
  \item \textbf{Operating point}: cấu hình đã chạy thật và pass thật.
  \item \textbf{Capacity decomposition}: không đánh đồng raw candidates với capacity thật.
\end{itemize}

\section{Kết Luận}

Hệ thống xử lý video theo hướng rất cụ thể:

\begin{center}
\texttt{H.264 bitstream}
\(\rightarrow\)
\texttt{IDR/CAVLC blocks}
\(\rightarrow\)
\texttt{T1 sign positions}
\(\rightarrow\)
\texttt{payload bits}
\(\rightarrow\)
\texttt{patched stego H.264}
\end{center}

Nó không giấu tin bằng cách sửa pixel. Nó giấu tin bằng cách sửa rất nhỏ dấu
của một số residual coefficients đã được chọn lọc trong bitstream H.264. Chính
những lớp lọc như patchability, reconstruction, quality guard, và operating
capacity mới làm hệ thống khác với một phương pháp nhúng đơn giản.

\end{document}
